% ============================================================
%  Chapter 3 -- Root and Ratio Tests, Power Series,
%               Summation by Parts, Absolute Convergence: 3.33-3.46
% ============================================================
\rsection{The Root and Ratio Tests}

\begin{signpost}[Two comparisons with the geometric series]
Neither test is a new idea. Both ask one question --- does $\Sigma a_n$ sit
below some geometric series $\Sigma\beta^n$ with $\beta < 1$? --- and differ
only in how they detect it: the root test reads the geometric rate off
$\sqrt[n]{|a_n|}$ directly, the ratio test off the quotient of consecutive
terms. This is why $\limsup$ was built: the rates need not converge, and
3.17(b) is exactly the tool that converts ``$\limsup < 1$'' into ``eventually
below a fixed $\beta < 1$.''

\smallskip
The relationship between the two is settled by 3.37: the root test is
strictly stronger, and Example 3.35 exhibits a series the ratio test cannot
touch. Use the ratio test for factorials and products, the root test for
theory --- it is the one the radius of convergence (3.39) is built on.
\end{signpost}

\ritem{3.33}{Theorem}
\emph{Given $\Sigma a_n$, put
$\alpha = \limsup_{n\to\infty} \sqrt[n]{|a_n|}$.}
\begin{enumerate}[label=(\alph*), leftmargin=2.2em, itemsep=1pt, topsep=3pt]
  \item \emph{If $\alpha < 1$, $\Sigma a_n$ converges;}
  \item \emph{if $\alpha > 1$, $\Sigma a_n$ diverges;}
  \item \emph{if $\alpha = 1$, the test gives no information.}
\end{enumerate}

\begin{proof}
If $\alpha < 1$, we can choose $\beta$ so that $\alpha < \beta < 1$, and an
integer $N$ such that
\[ \sqrt[n]{|a_n|} < \beta \]
for $n \ge N$ [by Theorem 3.17(b)]. That is, $n \ge N$ implies
\[ |a_n| < \beta^n. \]
Since $0 < \beta < 1$, $\Sigma\beta^n$ converges. Convergence of $\Sigma a_n$
follows now from the comparison test.\kn{Watch 3.17(b) do precisely its
promised job: $\beta > \alpha = s^*$ caps the whole tail. The room between
$\alpha$ and $1$ is essential --- the comparison is with $\Sigma\beta^n$ for a
\emph{fixed} $\beta < 1$, which is why $\alpha = 1$ proves nothing: no room
remains.}

If $\alpha > 1$, then, again by Theorem 3.17, there is a sequence $(n_k)$ such
that
\[ \sqrt[n_k]{|a_{n_k}|} \to \alpha. \]
Hence $|a_n| > 1$ for infinitely many values of $n$, so that the condition
$a_n \to 0$, necessary for convergence of $\Sigma a_n$, does not hold (Theorem
3.23).

To prove (c), we consider the series
\[ \sum \frac1n, \qquad \sum \frac{1}{n^2}. \]
For each of these series $\alpha = 1$, but the first diverges and the second
converges.\kn{Both $\alpha$'s are $1$ by 3.20(c): $\sqrt[n]{1/n} =
1/\sqrt[n]{n} \to 1$, and likewise for $1/n^2$. The $p$-series live exactly on
the geometric tests' blind spot --- their terms decay polynomially, slower
than every geometric rate, which is why 3.27/3.28 needed a different idea.}
\end{proof}

\ritem{3.34}{Theorem}
\emph{The series $\Sigma a_n$}
\begin{enumerate}[label=(\alph*), leftmargin=2.2em, itemsep=1pt, topsep=3pt]
  \item \emph{converges if
        $\displaystyle\limsup_{n\to\infty}
        \left|\frac{a_{n+1}}{a_n}\right| < 1$,}
  \item \emph{diverges if $\left|\dfrac{a_{n+1}}{a_n}\right| \ge 1$ for all
        $n \ge n_0$, where $n_0$ is some fixed integer.}
\end{enumerate}

\begin{proof}
If condition (a) holds, we can find $\beta < 1$, and an integer $N$, such that
\[ \left|\frac{a_{n+1}}{a_n}\right| < \beta \]
for $n \ge N$. In particular,
\begin{align*}
  |a_{N+1}| &< \beta|a_N|,\\
  |a_{N+2}| &< \beta|a_{N+1}| < \beta^2|a_N|,\\
  &\;\;\vdots\\
  |a_{N+p}| &< \beta^p|a_N|.
\end{align*}
That is,
\[ |a_n| < |a_N|\beta^{-N}\cdot\beta^n \]
for $n \ge N$, and (a) follows from the comparison test, since $\Sigma\beta^n$
converges.

If $|a_{n+1}| \ge |a_n|$ for $n \ge n_0$, it is easily seen that the condition
$a_n \to 0$ does not hold, and (b) follows.
\end{proof}

\emph{Note:} the knowledge that $\lim a_{n+1}/a_n = 1$ implies nothing about
the convergence of $\Sigma a_n$. The series $\sum 1/n$ and $\sum 1/n^2$
demonstrate this.

\begin{pitfall}[The divergence half of the ratio test is not a limsup]
Read (b) again: it demands $|a_{n+1}/a_n| \ge 1$ \emph{for all} $n \ge n_0$
--- a pointwise, eventually-everywhere condition. The tempting statement
``$\limsup |a_{n+1}/a_n| > 1$ implies divergence'' is \textbf{false}, and
Example 3.35(b) below is the counterexample: a convergent series with ratio
$\limsup$ equal to $2$. A few large ratios do no harm if they are followed by
crashes; only a ratio that \emph{never} again dips below $1$ pins the terms
away from zero.

Symmetrically, the convergence half really is a $\limsup$, and the
$\liminf$-version ``$\liminf |a_{n+1}/a_n| < 1$ implies convergence'' is false
too --- 3.35(a) has ratio $\liminf$ $0$ and could easily be built divergent.
The test is asymmetric because it detects geometric \emph{domination}, and
domination is one-sided.
\end{pitfall}

\ritem{3.35}{Examples}
\begin{enumerate}[label=(\alph*), leftmargin=2.2em, itemsep=4pt, topsep=4pt]
  \item Consider the series
        \[ \frac12 + \frac13 + \frac{1}{2^2} + \frac{1}{3^2} + \frac{1}{2^3}
           + \frac{1}{3^3} + \frac{1}{2^4} + \frac{1}{3^4} + \cdots, \]
        for which
        \begin{align*}
          \liminf_{n\to\infty} \frac{a_{n+1}}{a_n}
            &= \lim_{n\to\infty}\left(\frac23\right)^n = 0,\\
          \liminf_{n\to\infty} \sqrt[n]{a_n}
            &= \lim_{n\to\infty}\sqrt[2n]{\frac{1}{3^n}} = \frac{1}{\sqrt3},\\
          \limsup_{n\to\infty} \sqrt[n]{a_n}
            &= \lim_{n\to\infty}\sqrt[2n]{\frac{1}{2^n}} = \frac{1}{\sqrt2},\\
          \limsup_{n\to\infty} \frac{a_{n+1}}{a_n}
            &= \lim_{n\to\infty}\frac12\left(\frac32\right)^n = +\infty.
        \end{align*}
        The root test indicates convergence; the ratio test does not
        apply.\kn{The series interleaves two geometric series with different
        rates. The \emph{roots} feel only the rates, and both are safely below
        $1$; the \emph{ratios} feel the jumps between the two strands, which
        grow without bound. Interleaving is the standard way to break the
        ratio test.}
  \item The same is true for the series
        \[ \frac12 + 1 + \frac18 + \frac14 + \frac{1}{32} + \frac{1}{16}
           + \frac{1}{128} + \frac{1}{64} + \cdots, \]
        where
        \[ \liminf_{n\to\infty} \frac{a_{n+1}}{a_n} = \frac18, \qquad
           \limsup_{n\to\infty} \frac{a_{n+1}}{a_n} = 2, \]
        but
        \[ \lim_{n\to\infty} \sqrt[n]{a_n} = \frac12. \]
\end{enumerate}

\ritem{3.36}{Remarks}
The ratio test is frequently easier to apply than the root test, since it is
usually easier to compute ratios than $n$th roots. However, the root test has
wider scope. More precisely: Whenever the ratio test shows convergence, the
root test does too; whenever the root test is inconclusive, the ratio test is
too. This is a consequence of Theorem 3.37 and is illustrated by the above
examples.

Neither of the two tests is subtle with regard to divergence. Both deduce
divergence from the fact that $a_n$ does not tend to zero as $n \to \infty$.

\ritem{3.37}{Theorem}
\emph{For any sequence $(c_n)$ of positive numbers,}
\begin{align*}
  \liminf_{n\to\infty} \frac{c_{n+1}}{c_n}
    &\le \liminf_{n\to\infty} \sqrt[n]{c_n},\\
  \limsup_{n\to\infty} \sqrt[n]{c_n}
    &\le \limsup_{n\to\infty} \frac{c_{n+1}}{c_n}.
\end{align*}

\begin{proof}
We shall prove the second inequality; the proof of the first is quite similar.
Put
\[ \alpha = \limsup_{n\to\infty} \frac{c_{n+1}}{c_n}. \]
If $\alpha = +\infty$, there is nothing to prove. If $\alpha$ is finite, choose
$\beta > \alpha$. There is an integer $N$ such that
\[ \frac{c_{n+1}}{c_n} \le \beta \]
for $n \ge N$. In particular, for any $p > 0$,
\[ c_{N+k+1} \le \beta c_{N+k} \qquad (k = 0,1,\dots,p-1). \]
Multiplying these inequalities, we obtain
\[ c_{N+p} \le \beta^pc_N, \]
or
\[ c_n \le c_N\beta^{-N}\cdot\beta^n \qquad (n \ge N). \]
Hence
\[ \sqrt[n]{c_n} \le \sqrt[n]{c_N\beta^{-N}}\cdot\beta, \]
so that
\begin{equation}
  \limsup_{n\to\infty} \sqrt[n]{c_n} \le \beta, \tag{3.18}
\end{equation}
by Theorem 3.20(b). Since (3.18) is true for every $\beta > \alpha$, we have
\[ \limsup_{n\to\infty} \sqrt[n]{c_n} \le \alpha. \]\kn{The mechanism: a bound
on every \emph{ratio} compounds into a geometric bound on the terms
themselves, and taking $n$th roots flattens the constant $c_N\beta^{-N}$ into
$1$ by 3.20(b). Roots retain only the exponential rate; ratios carry every
local fluctuation. That is the entire content of ``root is stronger than
ratio,'' and it is also why $\lim c_{n+1}/c_n = L$ forces
$\lim\sqrt[n]{c_n} = L$ --- squeeze both displayed inequalities. Setting
$c_n = n!/n^n$ gives the useful byproduct $\sqrt[n]{n!}\sim n/e$.}
\end{proof}

\rsection{Power Series}

\begin{signpost}[Six chapters of consequences begin here]
Two items only --- a definition and one theorem --- but 3.39 is the
foundation on which Chapter 8 builds the exponential, the trigonometric
functions and eventually $e^{i\pi} = -1$. The theorem says the convergence
domain of $\Sigma c_nz^n$ is always a disc, and computes its radius by the
root test. That a single formula answers ``for which $z$?'' is the first sign
of how rigid power series are; the second, uniform convergence inside the
disc, arrives at 8.1.
\end{signpost}

\ritem{3.38}{Definition}
Given a sequence $(c_n)$ of complex numbers, the series
\begin{equation}
  \sum_{n=0}^{\infty} c_nz^n \tag{3.19}
\end{equation}
is called a \emph{power series}. The numbers $c_n$ are called the
\emph{coefficients} of the series; $z$ is a complex number.

In general, the series will converge or diverge, depending on the choice of
$z$. More specifically, with every power series there is associated a circle,
the \emph{circle of convergence}, such that (3.19) converges if $z$ is in the
interior of the circle and diverges if $z$ is in the exterior of the circle (to
cover all cases, we have to consider the plane as the interior of a circle with
infinite radius, and a point as a circle with radius zero). The behavior on the
circle of convergence is much more varied and cannot be described so simply.

\ritem{3.39}{Theorem}
\emph{Given the power series $\Sigma c_nz^n$, put}
\[ \alpha = \limsup_{n\to\infty} \sqrt[n]{|c_n|}, \qquad R = \frac{1}{\alpha}. \]
\emph{(If $\alpha = 0$, $R = +\infty$; if $\alpha = +\infty$, $R = 0$.) Then
$\Sigma c_nz^n$ converges if $|z| < R$ and diverges if $|z| > R$.}

\begin{proof}
Put $a_n = c_nz^n$ and apply the root test:
\[ \limsup_{n\to\infty} \sqrt[n]{|a_n|}
   = |z|\limsup_{n\to\infty} \sqrt[n]{|c_n|} = \frac{|z|}{R}. \]
\end{proof}

\emph{Note:} $R$ is called the \emph{radius of convergence} of
$\Sigma c_nz^n$.

\ritem{3.40}{Examples}
\begin{enumerate}[label=(\alph*), leftmargin=2.2em, itemsep=2pt, topsep=4pt]
  \item The series $\sum n^nz^n$ has $R = 0$.
  \item The series $\displaystyle\sum \frac{z^n}{n!}$ has $R = +\infty$. (In
        this case the ratio test is easier to apply than the root test.)
  \item The series $\sum z^n$ has $R = 1$. If $|z| = 1$, the series diverges,
        since $(z^n)$ does not tend to $0$ as $n \to \infty$.
  \item The series $\displaystyle\sum \frac{z^n}{n}$ has $R = 1$. It diverges
        if $z = 1$. It converges for all other $z$ with $|z| = 1$. (The last
        assertion will be proved in Theorem 3.44.)
  \item The series $\displaystyle\sum \frac{z^n}{n^2}$ has $R = 1$. It
        converges for all $z$ with $|z| = 1$, by the comparison test, since
        $|z^n/n^2| = 1/n^2$.
\end{enumerate}

\begin{center}
\begin{tikzpicture}[x=1mm, y=1mm]
  % the disc of convergence
  \draw[thin] (0,0) circle (17mm);
  \draw[axis] (-24,0) -- (26,0);
  \draw[axis] (0,-22) -- (0,24);
  \node[ptfill, minimum size=2.2pt] at (17,0) {};
  \node[mlbl, anchor=north west] at (16.5,-1.5) {$R$};
  \node[lbl] at (-7.5,7.5) {converges};
  \node[lbl] at (-7.2,4) {absolutely};
  \node[lbl, anchor=west] at (20,17) {diverges outside};
  \node[lbl, anchor=west] at (30,6) {on the circle:};
  \node[lbl, anchor=west] at (30,2.5) {all three behaviours occur ---};
  \node[lbl, anchor=west] at (30,-1) {3.40(c) nowhere, (d) everywhere};
  \node[lbl, anchor=west] at (30,-4.5) {but $z{=}1$, (e) everywhere};
  \draw[thin, -{Stealth[length=3pt]}] (30,5) -- (13,11.5);
\end{tikzpicture}
\end{center}
\figcap{Theorem 3.39: inside the circle everything converges (absolutely ---
the proof runs through the root test, hence 3.25), outside everything
diverges, and the circle itself carries no general rule.}

\rsection{Summation by Parts}

\begin{signpost}[The one tool for conditional convergence]
Rudin will observe at 3.46 that comparison, root and ratio only ever detect
\emph{absolute} convergence. This section supplies the single technique in the
book that goes further. The identity 3.41 trades $\Sigma a_nb_n$ for a sum
involving the partial sums $A_n$ --- profitable exactly when the $A_n$ are
bounded though $\Sigma a_n$ itself diverges, as with $a_n = (-1)^n$ or
$a_n = z^n$ on the unit circle. Theorem 3.42 packages the trade; 3.43
(alternating series) and 3.44 (boundary of the disc) are its two famous
corollaries.
\end{signpost}

\ritem{3.41}{Theorem}
\emph{Given two sequences $(a_n)$ and $(b_n)$, put}
\[ A_n = \sum_{k=0}^{n} a_k \]
\emph{if $n \ge 0$; put $A_{-1} = 0$. Then, if $0 \le p \le q$, we have}
\begin{equation}
  \sum_{n=p}^{q} a_nb_n = \sum_{n=p}^{q-1} A_n(b_n - b_{n+1})
  + A_qb_q - A_{p-1}b_p. \tag{3.20}
\end{equation}

\begin{proof}
\[ \sum_{n=p}^{q} a_nb_n = \sum_{n=p}^{q} (A_n - A_{n-1})b_n
   = \sum_{n=p}^{q} A_nb_n - \sum_{n=p-1}^{q-1} A_nb_{n+1}, \]
and the last expression on the right is clearly equal to the right side of
(3.20).
\end{proof}

Formula (3.20), the so-called ``partial summation formula,'' is useful in the
investigation of series of the form $\Sigma a_nb_n$, particularly when $(b_n)$
is monotonic. We shall now give applications.

\begin{deeper}[Integration by parts, one chapter early]
Rename the pieces and (3.20) is an old friend: $A_n$ is the ``antiderivative''
of $a_n$ (partial sums undo differencing, since $a_n = A_n - A_{n-1}$), and
$b_n - b_{n+1}$ is the ``derivative'' of $b_n$. The formula then reads
$\sum a\,b = -\sum A\,\Delta b + \text{boundary terms}$, which is
$\int u\,dv = uv - \int v\,du$ with sums for integrals. Chapter 6 proves the
integral version as 6.22, by essentially this computation applied to
Riemann--Stieltjes sums.

The strategic content is the same in both settings: move the burden from the
factor you cannot control onto the factor that varies slowly. Here $\Sigma a_n$
may diverge, but only the \emph{bounded} $A_n$ appear on the right, multiplied
by the increments of a monotone $b_n$ --- and those increments telescope.
\end{deeper}

\begin{center}
\begin{tikzpicture}[x=1mm, y=1mm]
  \fill[black!8] (10,0) rectangle (30,30);
  \fill[black!8] (30,0) rectangle (44,24);
  \fill[black!8] (44,0) rectangle (62,17);
  \fill[black!8] (62,0) rectangle (80,12);
  \fill[pattern=north east lines, pattern color=black!45]
    (44,0) rectangle (62,17);
  \fill[pattern=north west lines, pattern color=black!45]
    (10,24) rectangle (30,30);
  \draw[thin] (10,0) -- (10,30) -- (30,30) -- (30,24) -- (44,24) -- (44,17)
    -- (62,17) -- (62,12) -- (80,12) -- (80,0);
  \draw[thin] (30,0) -- (30,24);
  \draw[thin] (44,0) -- (44,17);
  \draw[thin] (62,0) -- (62,12);
  \draw[guide] (10,24) -- (30,24);
  \draw[guide] (10,17) -- (44,17);
  \draw[guide] (10,12) -- (62,12);
  \draw[axis] (4,0) -- (86,0);
  \foreach \x in {10,30,44,62,80}{\draw[thin] (\x,-1.6) -- (\x,0);}
  \node[mlbl, anchor=north] at (10,-2.2) {$A_{p-1}$};
  \node[mlbl, anchor=north] at (30,-2.2) {$A_p$};
  \node[mlbl, anchor=north] at (44,-2.2) {$A_{p+1}$};
  \node[mlbl, anchor=north] at (62,-2.2) {$A_{p+2}$};
  \node[mlbl, anchor=north] at (80,-2.2) {$A_q$};
  \node[mlbl, fill=white, inner sep=1pt] at (53,8) {$a_nb_n$};
  \node[mlbl, anchor=east] at (7.6,27)
    {$(A_n{-}A_{p-1})(b_n{-}b_{n+1})$};
  \node[mlbl, anchor=south] at (20,32) {$b_p$};
  \node[mlbl, anchor=west] at (82,6) {$b_q$};
\end{tikzpicture}
\end{center}
\figcap{Both sides of (3.20) measure this one staircase: columns of height
$b_n$ stand over the increments $a_n = A_n - A_{n-1}$. Sliced vertically
(hatched sample), the area is $\sum_p^q a_nb_n$. Sliced horizontally, the slab
between levels $b_{n+1}$ and $b_n$ has width $A_n - A_{p-1}$, and the base
slab, height $b_q$, runs the full width $A_q - A_{p-1}$; expanding the widths
telescopes all the $A_{p-1}$'s into the single boundary term $-A_{p-1}b_p$ of
(3.20). Drawn for $a_n > 0$ and decreasing positive $b_n$ --- exactly 3.42's
hypotheses; Rudin's two-line proof is the same bookkeeping with signs allowed.}

\ritem{3.42}{Theorem}
\emph{Suppose}
\begin{enumerate}[label=(\alph*), leftmargin=2.2em, itemsep=1pt, topsep=3pt]
  \item \emph{the partial sums $A_n$ of $\Sigma a_n$ form a bounded sequence;}
  \item \emph{$b_0 \ge b_1 \ge b_2 \ge \cdots$;}
  \item \emph{$\lim_{n\to\infty} b_n = 0$.}
\end{enumerate}
\emph{Then $\Sigma a_nb_n$ converges.}

\begin{proof}
Choose $M$ such that $|A_n| \le M$ for all $n$. Given $\varepsilon > 0$, there
is an integer $N$ such that $b_N \le \varepsilon/(2M)$. For
$N \le p \le q$, we have
\begin{align*}
  \left|\sum_{n=p}^{q} a_nb_n\right|
    &= \left|\sum_{n=p}^{q-1} A_n(b_n-b_{n+1}) + A_qb_q - A_{p-1}b_p\right|\\
    &\le M\left[\sum_{n=p}^{q-1} (b_n-b_{n+1}) + b_q + b_p\right]\\
    &= 2Mb_p \le 2Mb_N \le \varepsilon.
\end{align*}
Convergence now follows from the Cauchy criterion. We note that the first
inequality in the above chain depends of course on the fact that
$b_n - b_{n+1} \ge 0$.\kn{The telescope is the payoff: the increments
$b_n - b_{n+1}$ are nonnegative, so absolute values cost nothing, and their
sum collapses to $b_p - b_q$. All the oscillation of $\Sigma a_n$ has been
absorbed into the single constant $M$; what remains converges because $(b_n)$
does. Note the shape of the conclusion: convergence via 3.22, no sum ever
named.}
\end{proof}

\ritem{3.43}{Theorem}
\emph{Suppose}
\begin{enumerate}[label=(\alph*), leftmargin=2.2em, itemsep=1pt, topsep=3pt]
  \item \emph{$|c_1| \ge |c_2| \ge |c_3| \ge \cdots$;}
  \item \emph{$c_{2m-1} \ge 0$, $c_{2m} \le 0$ \quad $(m = 1,2,3,\dots)$;}
  \item \emph{$\lim_{n\to\infty} c_n = 0$.}
\end{enumerate}
\emph{Then $\Sigma c_n$ converges.}

Series for which (b) holds are called ``alternating series''; the theorem was
known to Leibnitz.

\begin{proof}
Apply Theorem 3.42 with $a_n = (-1)^{n+1}$, $b_n = |c_n|$.\kn{The partial sums
of $(-1)^{n+1}$ are $1,0,1,0,\dots$ --- bounded by $M = 1$ --- and the
$|c_n|$ decrease to $0$. The most famous instance:
$1 - \frac12 + \frac13 - \frac14 + \cdots$ converges --- its sum turns out
to be $\log 2$, a fact that needs Chapter 8 --- while its absolute version,
the harmonic series, diverges. The gap between those two facts is the subject of the whole
next section.}
\end{proof}

\begin{center}
\begin{tikzpicture}[x=1mm, y=1mm]
  \draw[axis] (-2,0) -- (96,0);
  \node[lbl, anchor=north] at (94,-2) {$n$};
  % partial sums of 1 - 1/2 + 1/3 - ... scaled: s=limit ~ .693 -> y=14
  \draw[guide] (0,14) -- (92,14);
  \node[mlbl, anchor=west] at (92.5,14) {$s$};
  % s1=1 -> 20.2, s2=.5 -> 10.1, s3=.833->16.8, s4=.583->11.8, s5=.783->15.8,
  % s6=.617->12.5, s7=.760->15.3, s8=.635->12.8, s9=.746->15.1, s10=.646->13.0
  \foreach \i/\x/\y in {1/6/20.2, 2/14/10.1, 3/22/16.8, 4/30/11.8, 5/38/15.8,
                        6/46/12.5, 7/54/15.3, 8/62/12.8, 9/70/15.1, 10/78/13.0}{
    \node[ptfill, minimum size=2.2pt] (s\i) at (\x,\y) {};}
  \foreach \i [evaluate=\i as \j using int(\i+1)] in {1,...,9}{
    \draw[thin] (s\i) -- (s\j);}
  \node[mlbl, anchor=south] at (6,21.5) {$s_1$};
  \node[mlbl, anchor=north] at (14,8.6) {$s_2$};
  \node[mlbl, anchor=south] at (22,18.2) {$s_3$};
  \node[lbl, anchor=west] at (30,22) {odd sums fall, even sums rise,};
  \node[lbl, anchor=west] at (30,18.5) {the limit is pinched between};
\end{tikzpicture}
\end{center}
\figcap{Alternating series: each swing overshoots the limit and is shorter
than the last. The even partial sums increase, the odd ones decrease, and
$|s - s_n| \le |c_{n+1}|$ --- the error is at most the first omitted term.}

\ritem{3.44}{Theorem}
\emph{Suppose the radius of convergence of $\Sigma c_nz^n$ is $1$, and suppose
$c_0 \ge c_1 \ge c_2 \ge \cdots$, $\lim_{n\to\infty} c_n = 0$. Then
$\Sigma c_nz^n$ converges at every point on the circle $|z| = 1$, except
possibly at $z = 1$.}

\begin{proof}
Put $a_n = z^n$, $b_n = c_n$. The hypotheses of Theorem 3.42 are then
satisfied, since
\[ |A_n| = \left|\sum_{m=0}^{n} z^m\right|
   = \left|\frac{1-z^{n+1}}{1-z}\right| \le \frac{2}{|1-z|}, \]
if $|z| = 1$ and $z \ne 1$.\kn{The bound degrades as $z$ approaches $1$ ---
$|1-z|$ in the denominator --- which is why $z = 1$ must be excepted; there
$A_n = n+1$ is unbounded and 3.42 has nothing to grip. This settles 3.40(d):
$\sum z^n/n$ converges everywhere on the unit circle except at $z=1$, where it
is the harmonic series. Alternating series are the case $z = -1$.}
\end{proof}

\rsection{Absolute Convergence}

\noindent
The series $\Sigma a_n$ is said to \emph{converge absolutely} if the series
$\Sigma|a_n|$ converges.

\ritem{3.45}{Theorem}
\emph{If $\Sigma a_n$ converges absolutely, then $\Sigma a_n$ converges.}

\begin{proof}
The assertion follows from the inequality
\[ \left|\sum_{k=n}^{m} a_k\right| \le \sum_{k=n}^{m} |a_k|, \]
plus the Cauchy criterion.
\end{proof}

\ritem{3.46}{Remarks}
For series of positive terms, absolute convergence is the same as convergence.

If $\Sigma a_n$ converges, but $\Sigma|a_n|$ diverges, we say that
$\Sigma a_n$ converges \emph{conditionally} [[nonabsolutely]].\kn{The doubled
brackets are the retyped edition's editorial mark, explained in its foreword:
Rudin's own word was ``nonabsolutely,'' silently modernised to
``conditionally'' throughout, with the original recorded once --- here --- at
the point of first use.} For instance, the series
\[ \sum_{n=1}^{\infty} \frac{(-1)^n}{n} \]
converges conditionally (Theorem 3.43).

The comparison test, as well as the root and ratio tests, is really a test for
absolute convergence, and therefore cannot give any information about
conditionally convergent series. Summation by parts can sometimes be used to
handle the latter. In particular, power series converge absolutely in the
interior of the circle of convergence.

We shall see that we may operate with absolutely convergent series very much as
with finite sums. We may multiply them term by term and we may change the order
in which the additions are carried out, without affecting the sum of the
series. But for conditionally convergent series this is no longer true, and
more care has to be taken when dealing with them.\kd{This paragraph is the
thesis statement for the rest of the chapter. ``Absolutely convergent'' will
turn out to mean ``behaves like a finite sum'' --- products work (3.50),
rearrangements are harmless (3.55) --- while ``conditionally convergent''
means the sum is an artifact of the \emph{order} of summation, so much so that
rearrangement can produce any sum whatever (3.54). Absolute convergence is not
a technicality; it is the boundary between arithmetic and analysis.}
